% ch04_generalized_spaces.tex — Generalized Metric Spaces
\chapter{Generalized Metric Spaces}\label{ch:generalized_spaces}

\section{Introduction}

The notion of metric space, as Fr\'echet defined it in 1906, rests on three simple axioms. But what happens if we relax one of them? If the triangle inequality only holds up to a multiplicative constant, we obtain a \emph{$b$-metric space}. If the distance from a point to itself is not necessarily zero, we enter the world of \emph{partial metric spaces}, invented by Steve Matthews in 1994 to model computational domains. Each generalisation arises from a concrete need---fractal analysis, program semantics, measure theory---and each raises the same question: does the fixed point theorem survive?

This chapter shows that the answer is yes, with subtle adaptations, and presents $b$-metric spaces,
partial metric spaces, and $G$-metric spaces, with their associated fixed point
theorems.

\section{$b$-Metric spaces}

\begin{definition}[$b$-metric space]
Let $X$ be a set and $s \geq 1$ a real number. A function
$d : X \times X \to [0, +\infty)$ is a \emph{$b$-metric} (with constant $s$)
if for all $x, y, z \in X$:
\begin{enumerate}
    \item $d(x, y) = 0 \Leftrightarrow x = y$;
    \item $d(x, y) = d(y, x)$;
    \item $d(x, z) \leq s \bigl[ d(x, y) + d(y, z) \bigr]$ \quad ($b$-triangle inequality).
\end{enumerate}
The pair $(X, d)$ is then called a \emph{$b$-metric space}.
\end{definition}

\begin{remark}
For $s = 1$, we recover the classical definition of a metric space. The constant
$s$ measures the ``defect'' of the triangle inequality.
\end{remark}

\begin{example}[Natural $b$-metric]
Let $X = \R$ and $d(x, y) = (x - y)^2$. Then $(X, d)$ is a $b$-metric space
with $s = 2$. Indeed, $(x - z)^2 = ((x-y) + (y-z))^2 \leq 2(x-y)^2 + 2(y-z)^2
= 2[d(x,y) + d(y,z)]$ by the inequality $(a + b)^2 \leq 2(a^2 + b^2)$.

However, $(X, d)$ is not a classical metric space since
$d(0, 2) = 4 > d(0, 1) + d(1, 2) = 1 + 1 = 2$.
\end{example}

\begin{example}[$\ell^p$ space with $0 < p < 1$]
The space $\ell^p$ for $0 < p < 1$ equipped with $d(x, y) = \sum_{n=1}^\infty
\abs{x_n - y_n}^p$ is a $b$-metric space with $s = 2^{1/p - 1}$, but it is not
a metric space.
\end{example}

\begin{definition}[Completeness in a $b$-metric space]
The notions of Cauchy sequence and completeness are defined as in the classical
case: $(x_n)$ is Cauchy if $d(x_m, x_n) \to 0$ as $m, n \to \infty$, and
$(X, d)$ is complete if every Cauchy sequence converges.
\end{definition}

\begin{warning}[title=\textbf{Difficulties specific to $b$-metrics}]
In a $b$-metric space:
\begin{itemize}
    \item The induced topology may not be metrizable in the classical sense;
    \item A $b$-metric is not necessarily continuous (the function
          $(x, y) \mapsto d(x, y)$ may not be continuous for the induced topology);
    \item Open balls may not be open sets in the induced topology.
\end{itemize}
\end{warning}

\subsection{Contraction principle in $b$-metric spaces}

\begin{theorem}[Banach for $b$-metrics]\label{thm:banach-b}
Let $(X, d)$ be a complete $b$-metric space with constant $s \geq 1$, and
$T : X \to X$ a mapping satisfying
\[
    d(Tx, Ty) \leq k \, d(x, y) \quad \forall x, y \in X
\]
with $0 \leq k < 1/s$. Then $T$ has a unique fixed point.
\end{theorem}

\begin{proof}
Let $x_0 \in X$ and $x_n = T^n(x_0)$. We have $d(x_{n+1}, x_n) \leq k^n d(x_1, x_0)$.
For $m > n$, applying the $b$-triangle inequality recursively:
\begin{align*}
    d(x_n, x_m) &\leq s \, d(x_n, x_{n+1}) + s^2 d(x_{n+1}, x_{n+2}) + \cdots + s^{m-n} d(x_{m-1}, x_m) \\
    &\leq \sum_{i=0}^{m-n-1} s^{i+1} k^{n+i} d(x_1, x_0)
    = s k^n \sum_{i=0}^{m-n-1} (sk)^i \, d(x_1, x_0).
\end{align*}
Since $sk < 1$, the geometric series converges and
$d(x_n, x_m) \leq \frac{sk^n}{1 - sk} d(x_1, x_0) \to 0$.
So $(x_n)$ is Cauchy and converges to $x^* \in X$.

Using: $d(x^*, Tx^*) \leq s[d(x^*, x_{n+1}) + d(Tx_n, Tx^*)]
\leq s[d(x^*, x_{n+1}) + k \, d(x_n, x^*)] \to 0$.

Uniqueness follows the same argument as in the classical case.
\end{proof}

\begin{remark}
The condition $k < 1/s$ is more restrictive than $k < 1$. For $s = 1$ (classical
metric space), we recover the usual condition $k < 1$.
\end{remark}

\begin{theorem}[Improved version, Czerwik 1998]
If $(X, d)$ is a complete $b$-metric space and $T : X \to X$ satisfies
$d(Tx, Ty) \leq k \, d(x, y)$ with $k < 1$ (without the restriction $k < 1/s$),
then $T$ has a unique fixed point, provided $d$ is continuous.
\end{theorem}

\section{Partial metric spaces}

\begin{definition}[Partial metric space]
A \emph{partial metric space} is a pair $(X, p)$ where
$p : X \times X \to [0, +\infty)$ satisfies for all $x, y, z \in X$:
\begin{enumerate}
    \item $x = y \Leftrightarrow p(x, x) = p(x, y) = p(y, y)$;
    \item $p(x, x) \leq p(x, y)$ \quad (small self-distance);
    \item $p(x, y) = p(y, x)$ \quad (symmetry);
    \item $p(x, z) \leq p(x, y) + p(y, z) - p(y, y)$ \quad (modified triangle inequality).
\end{enumerate}
\end{definition}

\begin{remark}
The essential difference from a classical metric is that $p(x, x)$ can be
strictly positive. This is motivated by theoretical computer science: $p(x, x)$
represents the ``size'' or ``cost'' of the information $x$.
\end{remark}

\begin{example}
On $X = [0, +\infty)$, $p(x, y) = \max(x, y)$ is a partial metric. We have
$p(x, x) = x \geq 0$, with equality only at $x = 0$.
\end{example}

\begin{theorem}[Matthews, 1994]
Let $(X, p)$ be a complete partial metric space and $T : X \to X$ a contraction,
i.e., $p(Tx, Ty) \leq k \, p(x, y)$ with $k \in [0, 1)$. Then $T$ has a unique
fixed point $x^*$, and $p(x^*, x^*) = 0$.
\end{theorem}

\begin{proof}
The partial metric $p$ induces a classical metric
$d_p(x, y) = 2p(x, y) - p(x, x) - p(y, y)$, and $T$ is also a contraction for
$d_p$. By Banach's theorem applied to $(X, d_p)$ (which is complete if $(X, p)$
is), $T$ has a unique fixed point $x^*$. Moreover,
$p(x^*, x^*) = p(Tx^*, Tx^*) \leq k \, p(x^*, x^*) \leq k^2 p(x^*, x^*) \leq \ldots$,
hence $p(x^*, x^*) = 0$.
\end{proof}

\section{$G$-Metric spaces}

\begin{definition}[$G$-metric space (Mustafa--Sims)]
A \emph{$G$-metric space} is a pair $(X, G)$ where
$G : X \times X \times X \to [0, +\infty)$ satisfies for all $x, y, z, a \in X$:
\begin{enumerate}
    \item $G(x, y, z) = 0 \Leftrightarrow x = y = z$;
    \item $0 < G(x, x, y)$ for $x \neq y$;
    \item $G(x, x, y) \leq G(x, y, z)$ for $z \neq y$;
    \item $G(x, y, z) = G(x, z, y) = G(y, z, x) = \ldots$ (symmetry in all variables);
    \item $G(x, y, z) \leq G(x, a, a) + G(a, y, z)$ (rectangle inequality).
\end{enumerate}
\end{definition}

\begin{example}
For any metric space $(X, d)$, one can define $G(x, y, z) = d(x, y) + d(y, z) + d(x, z)$
or $G(x, y, z) = \max\{d(x,y), d(y,z), d(x,z)\}$. Both define $G$-metrics.
\end{example}

\begin{theorem}[Mustafa--Sims, 2006]
Let $(X, G)$ be a complete $G$-metric space and $T : X \to X$ satisfying
\[
    G(Tx, Ty, Tz) \leq k \, G(x, y, z) \quad \forall x, y, z \in X
\]
with $k \in [0, 1)$. Then $T$ has a unique fixed point.
\end{theorem}

\begin{proof}
Fix $x_0$ and set $x_n = T^n(x_0)$. We have:
\[
    G(x_n, x_{n+1}, x_{n+1}) = G(Tx_{n-1}, Tx_n, Tx_n) \leq k \, G(x_{n-1}, x_n, x_n)
    \leq k^n \, G(x_0, x_1, x_1).
\]
For $m > n$:
\[
    G(x_n, x_m, x_m) \leq G(x_n, x_{n+1}, x_{n+1}) + G(x_{n+1}, x_m, x_m)
    \leq \sum_{i=n}^{m-1} k^i G(x_0, x_1, x_1) \leq \frac{k^n}{1-k} G(x_0, x_1, x_1).
\]
So $(x_n)$ is $G$-Cauchy and converges to $x^* \in X$.

$G(x^*, Tx^*, Tx^*) \leq G(x^*, x_{n+1}, x_{n+1}) + G(x_{n+1}, Tx^*, Tx^*)
\leq G(x^*, x_{n+1}, x_{n+1}) + k \, G(x_n, x^*, x^*) \to 0$. Hence $Tx^* = x^*$.
\end{proof}

\begin{remark}
Jleli and Samet (2012) showed that many fixed point results in $G$-metric
spaces can be deduced from classical results in ordinary metric spaces via the
associated metric $d_G(x, y) = G(x, y, y) + G(x, x, y)$. This observation has
sparked debate about the actual utility of the $G$-metric concept.
\end{remark}

\section{Rectangular metric spaces}

\begin{definition}[Rectangular metric space (Branciari)]
A rectangular metric space is a pair $(X, d)$ where
$d : X \times X \to [0, +\infty)$ satisfies axioms (1) and (2) of a metric
space, and where the triangle inequality is replaced by the \emph{rectangular
inequality}: for all distinct $x, y$ and all $u, v \in X \setminus \{x, y\}$
distinct,
\[
    d(x, y) \leq d(x, u) + d(u, v) + d(v, y).
\]
\end{definition}

\begin{theorem}[Branciari, 2000]
Let $(X, d)$ be a complete rectangular metric space and $T : X \to X$ a Banach
contraction ($d(Tx, Ty) \leq k \, d(x,y)$, $k < 1$). Then $T$ has a unique
fixed point.
\end{theorem}

\begin{warning}[title=\textbf{Subtleties}]
Branciari's original proof contained gaps, notably because in a rectangular
space, limits are not necessarily unique and the topology can be pathological.
Corrections were provided by Sarma et al.\ (2009).
\end{warning}

\section{Modular metric spaces}

\begin{definition}[Modular metric space]
Let $X$ be a set. A function $\omega : (0, \infty) \times X \times X \to [0, +\infty)$
is a \emph{metric modular} if for all $x, y, z \in X$ and all $\lambda, \mu > 0$:
\begin{enumerate}
    \item $\omega(\lambda, x, y) = 0$ for all $\lambda > 0$ $\Leftrightarrow$ $x = y$;
    \item $\omega(\lambda, x, y) = \omega(\lambda, y, x)$;
    \item $\omega(\lambda + \mu, x, z) \leq \omega(\lambda, x, y) + \omega(\mu, y, z)$.
\end{enumerate}
\end{definition}

\begin{theorem}
Let $X_\omega = \{x \in X : \omega(\lambda, x, x_0) < \infty \text{ for some } \lambda > 0\}$
be the associated modular metric space. If $X_\omega$ is $\omega$-complete and
$T : X_\omega \to X_\omega$ satisfies
$\omega(\lambda, Tx, Ty) \leq k \, \omega(\lambda, x, y)$ for all $\lambda > 0$
with $k < 1$, then $T$ has a unique fixed point.
\end{theorem}

\section{Comparison of generalizations}

\begin{keyformulas}[title=\textbf{Summary of generalized spaces}]
\begin{center}
\renewcommand{\arraystretch}{1.4}
\begin{tabular}{|l|c|c|}
\hline
\textbf{Space} & \textbf{Modification} & \textbf{FP condition} \\
\hline
Metric & -- & $k < 1$ \\
$b$-metric & $d(x,z) \leq s[d(x,y)+d(y,z)]$ & $k < 1/s$ or $k < 1$ \\
Partial & $p(x,x) \geq 0$ & $k < 1$ \\
$G$-metric & 3 variables & $k < 1$ \\
Rectangular & Rectangular ineq. & $k < 1$ \\
\hline
\end{tabular}
\end{center}
\end{keyformulas}

\section{Exercises}

\begin{exercise}
Show that $d(x,y) = \abs{x - y}^2$ defines a $b$-metric space on $\R$ with
$s = 2$, but not a metric space.
\end{exercise}

\begin{exercise}
Verify that $\ell^p$ for $0 < p < 1$ with $d(x,y) = \sum \abs{x_n - y_n}^p$ is
a complete $b$-metric space and compute the constant $s$.
\end{exercise}

\begin{exercise}
Let $(X, p)$ be a partial metric space. Show that
$d_p(x,y) = 2p(x,y) - p(x,x) - p(y,y)$ is a classical metric on $X$.
\end{exercise}

\begin{exercise}
Let $(X, d)$ be a metric space. Show that $G(x,y,z) = d(x,y) + d(y,z) + d(x,z)$
defines a $G$-metric space. What is the relationship between completeness of
$(X, d)$ and that of $(X, G)$?
\end{exercise}

\begin{exercise}
Construct a rectangular metric space that is not a metric space.
\textit{Hint: consider a four-point set with a carefully chosen metric.}
\end{exercise}

\begin{exercise}
Let $(X, d)$ be a complete $b$-metric space with constant $s = 2$ and
$T : X \to X$ with $d(Tx, Ty) \leq \frac{1}{3} d(x,y)$. Show that $T$ has a
unique fixed point and compute the a priori error estimate.
\end{exercise}
